\begin{lemma}[Claim 4.3(b3): exact edge effects of restoration]
Let \(\mathcal R\) be a restoration trace satisfying
\noderef{ExtensionBaseRestorationInvariant}, with initial coloring
\(\sigma_n\) and output \(\varrho_1\).  Put
\[
                         \mathcal S^*
   =\bigcup_{t\in L^*}S_t,
\]
using the exact stored pairs of the certified PE slots in
\(L^*=\noderef{ExtensionBaseRestorationIndices}\).  If an exact edge
identity \(f\) has
\(\varrho_1(f)\ne\sigma_n(f)\), then \(f\) is not internal to \(T_j\),
and there are colors \(c,d\in\mathcal S^*\) such that
\[
                     \sigma_n(f)=c,\qquad \varrho_1(f)=d.
\]
This is Claim 4.3(b3), with the two optional coloring values made
explicitly nonempty.
\end{lemma}
\begin{proof}
Write the trace as
\(\sigma^0,\ldots,\sigma^r\).  If the initial and final values on \(f\)
differ, at least one transition changes \(f\).  At every such transition,
the definition \noderef{TwoColorComponentSwap} says that the edge node
of \(f\) lies in the chosen component and that its value is exchanged
between the exact colors \(\gamma_{t_a}\) and \(\delta_{t_a}\).
The slot-data fields of
\noderef{ExtensionBaseRestorationTrace} identify those colors with the
certified gamma and delta fields.  The PE constructor in
\noderef{TashkinovSeriesStep} identifies its stored set as
\(S_{t_a}=\{\gamma_{t_a},\delta_{t_a}\}\).  Since \(t_a\in L^*\), every
value acquired in a transition belongs to \(\mathcal S^*\).

Choose the first transition which changes \(f\).  Before it, the value
is still \(\sigma_n(f)\), so the preceding paragraph writes that value
as \(\mathsf{some}(c)\) with \(c\in\mathcal S^*\).  Choose the last
transition which changes \(f\).  After it, the value is
\(\varrho_1(f)=\mathsf{some}(d)\) with
\(d\in\mathcal S^*\).  This proves the asserted color control, including
the fact that neither endpoint value is uncolored.

It remains to exclude an internal edge.  At a transition which changes
\(f\), the invariant says that the chosen component meets \(T_{t_a}\)
in exactly its support \(v_{t_a}\), and that \(T_j\) is an ordered
initial segment of \(T_{t_a}\).  If \(f\) were internal to \(T_j\), both
of its endpoint vertex nodes would lie in \(T_{t_a}\).  Each is adjacent
to the edge node of \(f\) in \noderef{TwoColorIncidenceGraph}, so both
would belong to the chosen component and hence both would equal
\(v_{t_a}\).  This would make \(f\) a loop, contrary to the loopless law
in \noderef{LooplessMultigraph}.  Thus \(f\) is not internal to \(T_j\),
which completes (b3) without identifying parallel edges.
\end{proof}
